% https://github.com/zongpingding/zTeX_bundle/tree/dev/Examples/example_x
% https://ask.latexstudio.net/ask/question/17926.html
% 
\documentclass[tikz,border=12pt]{standalone}
\usepackage[fontset=fandol]{ctex}
\usepackage{graphicx}
\usepackage{amsmath,amssymb}
\usepackage{forest}
\usepackage[export]{adjustbox}
\usetikzlibrary{arrows.meta}
\useforestlibrary{edges}
\forestset{
  root node/.style={draw,fill=red!85!black,text=black,rounded corners=3pt,align=center,font=\small,text width=4cm,inner xsep=5pt,inner ysep=3pt},
  category node/.style={draw,fill=blue!75!black,text=white,rounded corners=3pt,align=center,font=\small,inner xsep=5pt,inner ysep=3pt},
  concept node/.style={draw,fill=green!50!black,text=white,rounded corners=3pt,align=center,font=\small,inner xsep=5pt,inner ysep=3pt},
  note node/.style={draw,fill=white,text=black,rounded corners=3pt,align=left,font=\small,inner xsep=5pt,inner ysep=3pt}
}
\begin{document}
\begin{forest}
for tree={
  grow'=0,
  parent anchor=east,
  child anchor=west,
  anchor=west,
  forked edge,
  fork sep=15pt,
  l sep=25pt,
  s sep=5pt,
  if level=0{root node}{
    if level=1{category node}{
      if level=2{concept node}{note node}
    }
  }
}
[\textbf{平面点集}\\$(S)$ 为 $\mathbb{R}^n$ 上的点集，其补集\\$\mathbb{R}^n\setminus S$ 记作 $S^c$
  [点
    [内点
      [{存在 $\bar{x}$ 的一个 $\delta$ 邻域 $O(\bar{x},\delta)$ 完全落在 $S$ 中}]
      [{全体的点记作 $\operatorname{int}S$ 或 $S^\circ$}]
    ]
    [界点]
    [外点]
    [聚点
      [定义
        [{定义 1：$\bar{x}$ 的任一邻域都包含 $S$ 中无数个点.}]
        [{定义 2：$\bar{x}$ 的任一邻域都包含 $S$ 中的至少一个点.}]
      ]
      [{全体的点记作：$S^d$ 或 $S'$}]
      [{直观理解：也就是点 $\bar{x}$ 附近是不是聚集了很多 $S$ 中的点.}]
      [{内点一定是聚点，界点一定是聚点，有限集没有聚点.}]
    ]
    [孤立点]
  ]
  [点集
    [开集]
    [相对开集]
    [闭集]
    [相对闭集]
    [紧集
      [开覆盖
        [{如果 $\mathbb{R}^n$ 的一组开覆盖 $\{U_\alpha\}$ 满足 $\displaystyle\bigcup_\alpha U_\alpha \supset S$，那么称 $\{U_\alpha\}$ 是 $S$ 的一个开覆盖}]
        [{设 $E$ 是任何一个度量空间的子集，为每个 $p\in E$ 选择一个半径 $r_p$（半径可以是不同的），那么集族 $N_{r_p}(p)\mid p\in E$ 是 $E$ 的一个开覆盖}]
      ]
      [紧集]
      [Heine-Borel 定理
        [{$\mathbb{R}^n$ 上的点集 $S$ 是紧集的充分必要条件是它为有界闭集}]
        [{一个测试插入图片:\quad\includegraphics[valign=m,width=4cm]{example-image-duck}}]
      ]
    ]
  ]
  [区域
    [开域]
    [闭域]
  ]
  [基本定理
    [闭矩形套定理]
    [Cantor 闭区间套定理]
    [{凝聚定理 (Bolzano-Weierstrass 定理)：$\mathbb{R}^n$ 中的有界点列一定有收敛子列（有界无限点集一定有聚点）.}]
    [{Cauchy 定理：收敛点列 $\leftrightarrow$ 基本点列}]
    [{紧性定理 (Heine-Borel 定理)：一个基本测试}]
  ]
]
\end{forest}
\end{document}
